\section{The control flow language}
\label{sec:cfl}

\subsection{The forager, as a query}

Start with the forager's problem, stated the way a developer would state it
first. There is a table of specimens. There is a cheap test, $\chi$. The
question is which specimens pass the test.

\begin{lstlisting}[language=SQL,morekeywords={SELECT,FROM,WHERE}]
SELECT   spec
FROM     prey
WHERE    chi(spec)
\end{lstlisting}

Three clauses. \texttt{FROM} says where the data lives. \texttt{WHERE} says
which of it we want. \texttt{SELECT} says what shape to hand back. The shape is
so useful that it has been reinvented in every decade since it was
introduced~\cite{Codd70,ChamberlinBoyce74}.

\begin{center}
\small
\begin{tabular}{@{}ll@{}}
\toprule
\textbf{Dialect} & \textbf{The same query} \\
\midrule
SQL & \texttt{SELECT spec FROM prey WHERE chi(spec)} \\
LINQ & \texttt{from spec in prey where Chi(spec) select spec} \\
XQuery (FLWOR) & \texttt{for \$spec in \$prey where chi(\$spec) return \$spec} \\
Python & \texttt{[spec for spec in prey if chi(spec)]} \\
Haskell & \texttt{[ spec | spec <- prey, chi spec ]} \\
\bottomrule
\end{tabular}
\end{center}

The four descendants are LINQ~\cite{LINQ}, XQuery's
FLWOR~\cite{XQuery}, and the list comprehension in the form Wadler gave
it~\cite{Wadler87}. The persistence of the shape across five languages with
nothing else in common
is the fact worth taking seriously. It is not a syntax fashion. The comprehension
is what data access looks like when you write down honestly what you are doing:
naming a source, filtering it, and projecting the survivors.

\subsection{What LINQ was actually for}

There is a step in this history that is usually told as a story about
convenience and is better told as a story about expressiveness.

SQL is a complete language for asking. It is not a language for acting. A query
returns a result set to \emph{someone else} --- an application, in another
language, which then does whatever is to be done. The boundary between the two
is where the impedance mismatch lives, and every ORM ever written is an attempt
to plaster over it.

LINQ exists for precisely this reason: it embeds the query comprehension
\emph{into} a host language, so that a developer can at last express the
\texttt{DO}. That is the whole point. Once \texttt{from-where-select} is a
first-class expression of C\# or F\#, the result of a query is an ordinary value
in a language that can act on it, and the mismatch is gone --- for reads.
XQuery's FLWOR has the same ambition inside XML, and the same limit; XQuery
Update~\cite{XQueryUpdate} was bolted on afterwards precisely because the
\texttt{return} clause
cannot cause anything.

Rholang takes the desideratum LINQ identified and goes the rest of the way.
Rather than embedding the comprehension in a host language, it builds the
language \emph{around} the comprehension. The comprehension is not a query
expression that yields a value; it is the sole way anything happens at all.

\subsection{\texttt{SELECT-FROM-WHERE-DO}}

Here is the forager's query again, in rholang.

\begin{lstlisting}
for( @spec <- prey where Chi(spec) ) {
  Break!(spec)
}
\end{lstlisting}

Clause by clause, against the SQL:

\begin{center}
\small
\begin{tabular}{@{}lll@{}}
\toprule
\textbf{SQL} & \textbf{rholang} & \textbf{What changed} \\
\midrule
\texttt{FROM prey} & \texttt{<- prey} & the source is a \emph{channel}, not a table \\
\texttt{WHERE chi(spec)} & \texttt{where Chi(spec)} & nothing yet; \S\ref{sec:logic} widens it \\
\texttt{SELECT spec} & \texttt{@spec} & the pattern \emph{is} the projection \\
--- & \texttt{\{ Break!(spec) \}} & the \texttt{DO} \\
\bottomrule
\end{tabular}
\end{center}

Two of those rows deserve more than a line.

\paragraph{The pattern is the projection, and it is stronger than one.}
A \texttt{SELECT} list picks columns. A rholang pattern matches structure: it
can destructure nested terms, bind sub-terms to variables, and --- as
\S\ref{sec:ddl} explains --- it matches \emph{up to the equations of the theory
that types the channel}, so two spellings of the same value match the same
pattern without the developer arranging it. A pattern is a projection in the
same sense that a regular expression is a substring test.

\paragraph{The \texttt{DO} is a process, not a value.}
The continuation \texttt{\{ Break!(spec) \}} is not returned to a caller. It
runs. It may send on channels, spawn further comprehensions, or do nothing.
There is no host language to hand results back to, because there is no outside.

\subsection{The other half: the write}

A query language has one privileged verb. Rholang has two, and the second is the
reason the first behaves differently than it does in SQL:

\begin{lstlisting}
prey!( spec )
\end{lstlisting}

That is the write. It is a term, it composes in parallel with everything else
using \texttt{|}, and --- this is the part with consequences --- it is
\emph{addressed the same way the read is}. The reader names \texttt{prey}; the
writer names \texttt{prey}; neither names the other.

Being explicit about both halves is what distinguishes this from a query
language with side effects bolted on. In SQL, \texttt{INSERT} and
\texttt{SELECT} are different kinds of statement, mediated by an engine that
owns the table. Here they are two terms in one program, and the table is not
owned by anybody.

\subsection{The translation table}

The following is probably the most useful page of this note for a reader coming
from SQL. Rholang distinguishes three ways of reading, and the distinction is
exactly the distinction a database developer already makes between a destructive
read, a standing query, and a plain look.

\begin{center}
\small
\begin{tabular}{@{}L{0.36\textwidth}L{0.56\textwidth}@{}}
\toprule
\textbf{rholang} & \textbf{SQL reading} \\
\midrule
\texttt{x!(Q)} &
\texttt{INSERT} --- put a row on \texttt{x} \\[3pt]
\texttt{for(p <- x)P} &
\texttt{SELECT ... FOR UPDATE}, then \texttt{DELETE}, then act. The datum is
consumed by exactly one reader. \\[3pt]
\texttt{for(p <= x)P} &
a standing query, or a trigger: the reader persists and fires again on the
next matching write. \\[3pt]
\texttt{for(p <<- x)P} &
plain \texttt{SELECT} --- a peek. Reads without consuming, so the datum remains
for others. \\[3pt]
\texttt{where cond} &
\texttt{WHERE} \\[3pt]
\texttt{\&} &
a join --- and simultaneously a distributed atomic commit across the joined
channels. \\[3pt]
\texttt{;} &
sequence: this group of receipts, then that one. \\[3pt]
the pattern &
the projection, structural rather than columnar, and matching up to the
theory's equations. \\
\bottomrule
\end{tabular}
\end{center}

The three arrows are worth committing to memory, because choosing among them is
the main design decision in a rholang program and the choice is invisible in the
SQL vocabulary that motivates it. A linear receipt \texttt{<-} is the default.
It consumes. Persistence \texttt{<=} and peeking \texttt{<<-} are the deviations,
and each gives something up: a persistent receipt gives up the guarantee that a
datum is handled once; a peek gives up the guarantee that reading it means
nobody else will.

\begin{remark}[Perception must not be able to act]
\label{rem:peek}
The difference between \texttt{<-} and \texttt{<<-} is not efficiency. It is
whether looking changes the world. A forager that reads the environment with
\texttt{<-} consumes what it observes; one that reads with \texttt{<<-} observes
without disturbing. In the running example this matters immediately: the
forager's cheap test $\chi$ must be a peek, or the act of assessing a specimen
would itself claim the specimen. Getting this wrong produces a system that
cannot observe anything without changing it, which is a real bug and not a
philosophical one.
\end{remark}

\subsection{Joins, groups, and where the \texttt{where} goes}

The full form of the comprehension is a sequence of groups, and each group is a
set of receipts that must occur \emph{together}:

\begin{lstlisting}
for( ptrn_11 <- chan_11 & ... & ptrn_1n <- chan_1n  where cond_1 ;
     ...                                                         ;
     ptrn_k1 <- chan_k1 & ... & ptrn_km <- chan_km  where cond_k ) P
\end{lstlisting}

Inside a group, \texttt{\&} means simultaneity: all of those receipts happen at
once or none of them do. Between groups, \texttt{;} means sequence. Each group
carries its own \texttt{where}, and \texttt{cond\_j} may mention any variable
bound by a pattern in group $i$ for $i \leq j$ --- so the guard on a later group
can test what an earlier group received. The comprehension commits only when
every pattern matches and every condition holds.

A SQL developer will recognise \texttt{\&} as a join and should immediately
notice the difference. A SQL join is over static extents: the rows are all
already there, and the join is a computation over them. A rholang join is over
\emph{concurrent writers}: the rows are being produced by other running
programs, and the join is a rendezvous. It succeeds when a matching datum is
available on every joined channel at once, and until then it waits.

\begin{aside}[A note for readers of the papers]
Several of the research notes in this line write a single trailing
\texttt{where} at the end of the whole comprehension. That is an abbreviation.
The implemented form is the per-group one above, and it is the better one to
learn from, because it makes visible exactly which bindings are in scope for
which condition.
\end{aside}

\subsection{The witnessed transaction}
\label{sec:transaction}

Now the payoff for having made both halves explicit.

When a reader and a writer meet on a channel, one rule fires. In the MeTTaIL
presentation of the calculus~\cite{mettailrust} it is written like this --- and this is a real
excerpt from a language definition, not pseudocode:

\begin{lstlisting}[language=mettail]
(PPar {(PInput n ^x.p), (POutput n q), ...rest})
    ~> (PPar {(subst ^x.p (NQuote q)), ...rest});
\end{lstlisting}

In words: a parallel composition containing a receipt on \texttt{n} and a send
on \texttt{n} rewrites to the same composition with both of them replaced by the
receipt's body, with the sent term substituted for the bound variable. Everything
else, \texttt{...rest}, is untouched.

The substitution is the point. After the step, the running term \emph{contains
the evidence} that this reader met this writer over this datum: the datum is
sitting in the continuation, in the position the pattern named. Nothing external
recorded the event. The event left a mark on the program.

\begin{claim}[Transactions are a property of the term]
\label{claim:witness}
A communication is atomic, and its atomicity requires no support from the
runtime. The condition is evaluated before the step commits; if it fails,
nothing is consumed, so there is no partial state and nothing to compensate. If
it succeeds, exactly one reader and one writer are consumed and the substitution
records that they were the ones. There is no state in which the read happened
and the write did not.
\end{claim}

This is not a subtle argument, and the implementation makes it concrete. The
comment at the head of the shipped example~\cite{f1r3node} says it plainly: a
failing guard
leaves the message in the tuple space, exactly as a non-matching pattern does. A
failed transaction is not rolled back; it never began.

\begin{lstlisting}
// Receive-side `where` guard: only commit when both the spatial pattern
// matches AND the guard evaluates to true. Failing the guard leaves the
// message in the tuple space -- same as a non-matching pattern.
new chan, stdout(`rho:io:stdout`) in {
    chan!(-3) |
    chan!(0)  |
    chan!(7)  |
    for (@x <- chan where x > 0) {
        stdout!(("kept", x))
    }
}
\end{lstlisting}

Three consequences a developer will care about.

\paragraph{Multi-party commit is a language construct.}
A join \texttt{\&} across $m$ channels is one transaction over $m$ resources. It
either takes a datum from each or takes none. The protocol machinery that would
normally implement this --- prepare, vote, commit, the coordinator, the timeout
--- is not written, because the rule that fires is the multi-channel one.

\paragraph{The witness is where cost attaches.}
Because there is an identified event, there is something to charge for. This is
what makes the metering of \S\ref{sec:cost} possible without instrumenting an
interpreter: the unit of billing is the same unit as the unit of commitment.

\paragraph{Candidates.}
For a join across several channels with several available data, there may be
more than one way to assign data to patterns. Each such assignment is a
\emph{candidate}. Ordinarily one of them happens and the others do not, and a
developer never thinks about it. Keep the word anyway: candidates are the thread
that runs through the rest of this note. In \S\ref{sec:logic} a condition is
evaluated per candidate; in \S\ref{sec:graded} a candidate carries a weight
rather than a verdict, and everything interesting follows from that.

\subsection{The forager, in full}

Enough machinery is now on the table to write the running example properly. The
forager has a belief, held as data on a channel. It races two courses of action:
open the specimen, or ignore it and graze a low-yield but safe alternative.

\begin{lstlisting}
new bel, prey, wall in {
  bel!( (1, 0.5) ) |                    // (n, s): one unit of prior evidence

  // the two branches race for the belief and the specimen
  for( @(n,s) <- bel ; @spec <- prey  where  Chi(spec) ) {
    Break!(spec) | bel!( revise(n, s, spec) )
  } |
  for( @(n,s) <- bel ; @spec <- prey  where  g ) {
    Graze!(*wall) | bel!( (n, s) )
  }
}
\end{lstlisting}

Read it as a developer. Two comprehensions compete for the same two resources:
the belief and the next specimen. Whichever fires consumes both, does its work,
and puts an updated belief back. Because the belief is consumed, the two
branches cannot both act on the same specimen --- mutual exclusion is not
arranged, it is what a linear receipt means.

Three questions are now unavoidable, and they are the agenda for the rest of
this note.

\begin{enumerate}[leftmargin=2.2em]
\item \texttt{Chi(spec)} is a condition about the \emph{structure} of a
  specimen. Where does the vocabulary for saying that come from?
  (\S\ref{sec:ddl}.)
\item What else could a condition say, beyond structure --- and what does saying
  more cost? (\S\ref{sec:logic}.)
\item \texttt{g} in the second branch is not a Boolean. It is a number. What is
  a numeric \texttt{where} clause, and what does a program mean when its
  conditions stop answering yes or no? (\S\ref{sec:graded}.)
\end{enumerate}
